% !TeX encoding = UTF-8 !TeX spellcheck = en_US
\documentclass[11pt, a4paper]{article}

% The title of the current document to be produced.
\newcommand{\doctitle}{Assignment \#3}
\newcommand{\assignments}{\bluetext{\textbf{Assignment:}}}
%
\setlength{\unitlength}{1in}
\renewcommand{\arraystretch}{2}
\setlength{\columnsep}{2cm}

\newcommand\dunderline[3][-1pt]{{%
  \sbox0{#3}%
  \ooalign{\copy0\cr\rule[\dimexpr#1-#2\relax]{\wd0}{#2}}}}

\usepackage[sfdefault]{atkinson}
\usepackage[T1]{fontenc}

\def\code#1{\bluetext{\texttt{#1}}}


%------------------------------------------------------------
% Import commands for both teacher and course information.  | NOTE: Change your
% teacher and course info in these files. |
%------>------>------>------>------>------>------>------>-->|
%
%------------------------------------------------------------
%-- Import packages and custom command definitons.          |
%------>------>------>------>------>------>------>------>-->|
%% ==================================
%% ===== Teacher info             ===
%% ==================================
%%
\newcommand{\instructor}{Katelyn Breivik}
\newcommand{\office}{Wean 8319}
\newcommand{\hours}{3pm on Mondays}
\newcommand{\phone}{}
\newcommand{\college}{}
\newcommand{\email}{kbreivik@andrew.cmu.edu}
\newcommand{\faculty}{}
\newcommand{\department}{Department of Physics}
                              %|
%% ==================================
%% ===== Course-specific commands ===
%% ==================================

%- Instructions: change course info here. 
\newcommand{\semester}{Fall 2023}

\newcommand{\ponderation}{2-4-3 (Theory-Lab-Homework)}
\newcommand{\coursetitle}{Astrophysics of Stars and the Galaxy}
\newcommand{\coursenumber}{[33-467]}
\newcommand{\prerequisite}{33-224, 33-234}
 
\input{../includes/packages-imports}                          %|  
\input{../includes/custom-commands}   
%
%---> Genereate & inject metadata                           |
\input{../includes/hyperef.doc.info}                          %|
%------------------------------------------------------------

\topmargin -70pt
\begin{document} 

\normalfont

%-------------------------------------------------------------
%-- Make the header of the document                          |
%------>------>------>------>------>------>------>------>--> |
%---------------------------------------------------------------------
%-- The following produces the document header including the title.  |
%---------------------------------------------------------------------

\noindent % <-- need to have this first.

\hfill	
\begin{minipage}{0.60\textwidth}%
	\raggedleft%
	{\Large \textsc{\doctitle}\par}
	\doublerule 
\end{minipage}%
\vspace{0.9cm}

  
%-------------------------------------------------------------
%-- Problem # 1                                              |
%------>------>------>------>------>------>------>------>--> |
\noindent \dunderline{1.0pt}{\textbf{Problem \#1: The $\rho$ - $T$ plane (22 points)}}\\

\noindent \textbf{Problem \#1a: Setting boundaries (10 points)}\\
Consider a gas of ionized hydrogen. Using the relations we derived in class find
the boundary in the $\rho$ - $T$ plane where:
\vspace{4pt}

\noindent ($i$) radiation pressure dominates the equation of state of the gas (5 points)
\vspace{4pt}

\noindent ($ii$) the ideal gas pressure dominates the equation of state of the
gas (5 points)\\

\noindent \textbf{Problem \#1b: Big stars and little stars (12 points)}\\

\noindent ($i$) Using the mass - radius relation calculate the average density
of a $10\,\rm{M_{\odot}}$ and $150\,\rm{M_{\odot}}$ star on the main sequence. (4 points)
\vspace{4pt}

\noindent ($ii$) Using the virial theorem as described in chapter two to
calculate the average temperature of the $10\,\rm{M_{\odot}}$ and
$150\,\rm{M_{\odot}}$ stars. (4 points)
\vspace{4pt}

\noindent ($iii$) Does radiation pressure or ideal gass pressure dominate for
each star based on their average densities and temperatures? (4 points)

\vspace{10pt}
\noindent \dunderline{1.0pt}{\textbf{Problem \#2: Writing an abstract from
preliminary work (20 points)}}\\
\noindent Working as a researcher, you'll encounter situations where you need to
write an abstract for a presentation on a research project that isn't finished
yet. This can be kind of a scary thing to do if you aren't quite sure what you
might actually end up finding; but... deadlines are deadlines! \\

\noindent For this problem, you should try to formulate as clear of an abstract
as you can based on your current progress in project 1. It is OK if you don't
have results yet! But hopefully you have a rough idea of what you hope to
accomplish and you can articulate why you think the topic is interesting. Every
good abstract starts with some context about your problem, then gives some
details for your methods, then hints that there could be interesting results.
The average length is between 5 and 7 sentences (ish), so this shouldn't take
you too long to write! \\

\noindent \emph{\large Hint!} If you need some inspiration, take a look at the
abstracts of the papers you searched for on ads in homework 2.

\end{document} 
